\section{Observabilidad distribuida}
\label{sec:observabilidad}

La observabilidad de un sistema distribuido se sustenta en tres señales complementarias:
métricas, registros estructurados y trazas correlacionadas por un identificador
común~\cite{sigelman2010dapper,beyer2016sre}. El estándar OpenTelemetry~\cite{opentelemetry2024spec}
unifica su modelo de datos y su exportación, permitiendo que distintos proveedores de \emph{backend}
consuman las mismas señales sin instrumentación duplicada.

\subsection{Métricas}

Se reutilizan las métricas de CockroachDB ya expuestas desde la Entrega 3
(\texttt{crdb\_query\_duration\_seconds}, \texttt{crdb\_transaction\_retries\_total},
\texttt{crdb\_pool\_active\_connections}) y se agregan cuatro métricas nuevas de aplicación,
recolectadas por Micrometer y expuestas en \texttt{/actuator/prometheus} de cada servicio Java (y
su equivalente en \texttt{api-gateway} vía un \emph{filter} reactivo):

\begin{itemize}[nosep]
  \item \texttt{http\_requests\_total} --- contador de peticiones, etiquetado por ruta, método y
        código de estado.
  \item \texttt{http\_request\_duration\_seconds} --- histograma de latencia, mismas etiquetas.
  \item \texttt{app\_active\_sessions} --- \emph{gauge} de sesiones activas reales, respaldado por
        el conteo de \emph{refresh tokens} no revocados y no expirados en \texttt{auth-service}.
  \item \texttt{app\_business\_events\_total} --- contador de eventos de negocio; el primer y
        único evento real que el sistema dispara por sí mismo es el escalado de un ticket
        (patrón Observer, Sección~\ref{sec:hexagonal-e4}).
\end{itemize}

En las rutas de \texttt{api-gateway} y \texttt{ticket-service}, los identificadores de recurso
(UUID) se normalizan a \texttt{\{id\}} en la etiqueta de ruta antes de registrar la métrica, para
evitar una explosión de cardinalidad de series temporales distintas por cada ticket individual.

\subsection{Registros estructurados}

Los cuatro servicios Java emiten registros en JSON estructurado con Logback
(\texttt{LogstashEncoder}), incluyendo el campo \texttt{trace\_id}. El agente Java de
OpenTelemetry, adjuntado en tiempo de ejecución vía \texttt{-javaagent}, escribe automáticamente
\texttt{trace\_id} y \texttt{span\_id} en el contexto de diagnóstico (MDC) de cada hilo que
procesa una petición instrumentada —la configuración de Logback solo necesita leer esas dos
claves, sin ningún cambio en el código de la aplicación.

\subsection{Trazas distribuidas}

El mismo agente Java exporta trazas vía OTLP a un \textbf{OpenTelemetry Collector}, que las
reenvía a \textbf{Grafana Tempo} para su almacenamiento y consulta. Una traza típica de una
petición \texttt{GET /api/v1/tickets} atraviesa: el span de recepción en \texttt{api-gateway}, el
span de la llamada HTTP saliente a \texttt{auth-service /validate}, el span del \emph{handler} de
aplicación en \texttt{ticket-service}, y el span de la consulta JDBC a CockroachDB —los cuatro
correlacionados por el mismo \texttt{trace\_id} que aparece en los registros de cada servicio
involucrado.

\textbf{Limitación conocida}: la aplicación móvil no incluye un SDK de OpenTelemetry propio, así
que la traza visible de una petición originada en el cliente móvil comienza en
\texttt{api-gateway}, no dentro del proceso del teléfono. Se documenta esta limitación
explícitamente en vez de presentar una instrumentación de cliente que no existe.

\subsection{Dashboard operativo}

El dashboard de Grafana, versionado como JSON reproducible en
\texttt{ops/grafana/pfc-dashboard.json}, incluye los seis paneles obligatorios: tasa de
peticiones, latencia p50/p95/p99, tasa de errores 5xx, disponibilidad Raft del cluster
CockroachDB (nodos vivos, rangos no disponibles), uso de recursos por contenedor (vía cAdvisor,
única forma práctica de obtener CPU/memoria por contenedor ya que los \texttt{/actuator/prometheus}
de cada servicio solo reportan métricas de su propia JVM), y latencia extremo a extremo, proxied
a través de las métricas de \texttt{api-gateway} por ser el único punto por el que pasan
absolutamente todas las peticiones de ambos clientes.
